Provjera sustava provedena je u tri koraka rastuće složenosti, gdje se u svakom sljedećem koraku mijenja točno jedna varijabla u odnosu na prethodni: prvo provjera upravljačkog sloja u njegovoj izvornoj konfiguraciji, zatim provjera prilagodbe za kameru okrenutu prema gore uz i dalje jednostavan i pouzdan cilj (ArUco oznaka), te naposljetku provjera cijelog sustava s cjevovodom za detekciju grane umjesto ArUco oznake. Ciljna točka u sva tri koraka do upravljačkog sloja stiže istim putem, UDP mostom opisanim u Odjeljku~\ref{subsec:udp-most}, tako da ta poveznica nije varijabla koja se mijenja između eksperimenata.

\section{Eksperiment 1: slijetanje na ArUco oznaku (kamera prema dolje)}
\label{sec:test1}

Cilj prvog eksperimenta jest provjeriti upravljački sloj (\texttt{ibvs\_perching}) u njegovoj izvornoj konfiguraciji, prije bilo kakve prilagodbe za scenarij prijanjanja na granu, te ugoditi odziv letjelice (parametre regulatora poravnanja) na ovoj jednostavnijoj konfiguraciji prije prelaska na složenije scenarije u drugom i trećem eksperimentu. Ciljnu točku i u ovom eksperimentu određuje isti cjevovod za detekciju opisan u poglavlju~\ref{ch:percepcija}, koji se koristi i u preostala dva eksperimenta, a ne modul za detekciju ArUco oznake ugrađen u paket \texttt{ibvs\_perching}. Razlika je u tome što cjevovod ovdje detektira ArUco oznaku umjesto grane, uz isključene filtre za ocjenjivanje kandidata koji su specifični za oblik i orijentaciju grane, dok ostatak cjevovoda, uključujući KLT praćenje i UDP most, ostaje nepromijenjen u sva tri eksperimenta.

\subsection{Postav}
\label{subsec:test1-postav}
ArUco oznaka postavljena je na tlo, a letjelica leti iznad nje s kamerom okrenutom prema dolje, izvorna konfiguracija paketa \texttt{ibvs\_perching}. Letjelica slijeće izravno na oznaku.

\section{Eksperiment 2: praćenje ArUco oznake (kamera prema gore)}
\label{sec:test2}

Cilj drugog eksperimenta jest izolirano provjeriti prilagodbu upravljačkog sloja za kameru okrenutu prema gore (Odjeljak~\ref{subsec:preslikavanje-osi}), zadržavajući ArUco oznaku kao jednostavan i pouzdan izvor ciljne točke. Time se promjena geometrije (kamera prema gore, cilj iznad letjelice) provjerava odvojeno od uvođenja cjevovoda za detekciju grane, koji je uveden tek u trećem eksperimentu. Budući da je kamera na letjelici radi ovog eksperimenta fizički zarotirana za 180 stupnjeva u odnosu na prvi eksperiment, cilj je ujedno i ispraviti preslikavanje osi tako da odgovara ovoj novoj orijentaciji kamere.

\subsection{Postav}
\label{subsec:test2-postav}
ArUco oznaka postavljena je na povišeno mjesto, a kamera na letjelici okrenuta je prema gore. Letjelica nastoji poravnati se s oznakom i prići joj odozdo.

\section{Eksperiment 3: prijanjanje na granu}
\label{sec:test3}

Cilj trećeg i glavnog eksperimenta jest provjeriti cijeli sustav opisan u ovom radu u konačnoj konfiguraciji: ArUco oznaka je uklonjena, a ciljnu točku određuje cjevovod za detekciju grane (poglavlje~\ref{ch:percepcija}), detekcija modelom, potom praćenje KLT algoritmom, čiji se izlaz UDP mostom prenosi upravljačkom sloju prilagođenom kameri okrenutoj prema gore (poglavlje~\ref{ch:upravljanje}).

\subsection{Postav}
\label{subsec:test3-postav}
Letjelica polazi pozicionirana neposredno ispod grane, s kamerom okrenutom prema gore. Potrebno je uzeti u obzir prostorni razmak između kamere i mehanizma za prijanjanje (hvataljke), budući da se ciljna točka određuje u odnosu na kameru, dok manevar prijanjanja mora dovesti granu u doseg hvataljke, a ne samo u središte slike.
